% Source: https://artofproblemsolving.com/community/c3h3181940_2022_abmc_online_october_contest__actonboxborough_math_competition
% Stable identity: aops-topic-3181940-q13
Write $F(t)=At^2+Bt+C$ with $A\ne0$, and put $U=x^2-x,V=2F(x)-1$. The identity factors as $(U-V)(A(U+V)+B)=0$. Since real polynomials have no zero divisors, one factor is identically zero. If $U=V$, coefficient comparison gives $F(x)=\dfrac{x^2-x+1}{2}$. If $A(U+V)+B=0$, it gives $A=-\dfrac{1}{2},B=\dfrac{1}{2},C=1$, so $F(x)=\dfrac{-x^2+x}{2}+1$. The two values at 2022 sum to $\dfrac{3}{2}$, hence $\boxed{a+b=5}$.
